\section{Background}\label{sec:background}
In the following, we illustrate the C-Learner in the context of the average treatment effect (ATE). We discuss extensions to other estimands in  \cref{sec:other_estimands}.

\subsection{Average Treatment Effect and Missing Outcomes}
\label{sec:ate}

\looseness=-1 Consider binary treatments (actions) $A\in\{0,1\}$, covariates $X\in \mathbb R^d$, and potential outcomes $Y(1), Y(0)\in \R$ under treatment ($A=1$) and control ($A=0$), respectively. 
The key difficulty in causal inference is that we do not observe counterfactuals:  
we only observe the outcome $Y := Y(A)$ corresponding to the observed binary treatment. %
Let $Z:=(X,A,Y)$. Based on i.i.d. observations $Z \sim P$, our goal is to estimate the average treatment effect $\psi(P) = P[Y(1) - Y(0)]:=\int [Y(1) - Y(0)]dP$. Here $P$ denotes both the joint probability measure and the expectation operator under this measure. %
We require standard identification conditions that make this goal feasible: 
(i) $Y=Y(A)$ (SUTVA), (ii) $(Y(1),Y(0))\perp A \mid X$ (ignorability), and (iii) for some $\eta>0$, $\;\eta\leq P(A=1\mid X)\leq 1-\eta$ a.s. (overlap) \citep{ding2023course}. 


To simplify our exposition, we assume $Y(0):=0$ throughout so that we are estimating the mean of a censored outcome (censored when $A=0$). 
This setting is also known as mean missing outcome \citep{Kennedy22}, which is the focus of some of our experiments. In this case, we can write the ATE as a functional of the joint measure $P$, as
\begin{equation}
    \label{eqn:ate}
\psi(P) := P[Y(1)]=P\left[P[Y\mid A=1, X]\right].
\end{equation}
The outcome model $\mu(X):=P[Y\mid A=1, X]$ (shorthand for $\mu(A,X):=P[Y\mid A, X]$ with $\mu(0,X):=0$)
and propensity (treatment) model $\pi(X) := P(A=1\mid X)$ are key nuisance parameters; notably, they are high-dimensional in contrast to the single-dimensional ATE~\eqref{eqn:ate}. 
Note that we can write $\psi(P)=P[\mu(1,X)-\mu(0,X)]=P[\mu(1,X)]=P[\mu(X)]$ in this setting. 
The corresponding nuisance estimators $\what \mu(X)$ and $\what \pi(X)$ can be implemented as ML models that are trained on held-out data; we will discuss data splits in \Cref{sec:data_splitting}. %

\subsection{First-Order Correction and Asymptotic Optimality} 
\label{sec:optimality_background}
By analyzing the error from blindly trusting an ML-based estimate of the nuisance parameters to estimate the ATE~\eqref{eqn:ate}, we can develop better estimation methods. %
We sketch intuition here and leave a more rigorous treatment to \cref{sec:theory_text}. 
\ifdefined\pnas
    Consider a statistical functional $Q \mapsto \psi(Q)$, with $\varphi$ its canonical gradient (a.k.a. efficient influence function) 
    with respect to $Q$; 
    without loss of generality, $Q[\varphi(Z; Q) | X] =0$ for all $Q$. 
    Then, a (distributional) first-order Taylor expansion gives
    \begin{align}
        \psi(\what{P}) 
        -  \psi(P)
        & = \int \varphi(Z; \what{P}) d (\what{P} - P)  + R_2(\what{P}, P) \nonumber \\
        & = - P [ \varphi(Z; \what{P}) ] 
        + R_2(\what{P}, P)
        \label{eqn:taylor}
    \end{align}
    where $R_2$ is a second-order remainder term.
    Conclude that the first-order error term of the plug-in approach $\psi(\what P)=\what P[\what{\mu}(X)]$ is given by $-P [ \varphi(Z; \what{P}) ]$.

\else
    We begin by noting that the joint distribution over the observed data can be decomposed as $P= P_{Y, A | X} \times P_X$.
    Since the marginal $P_X$ can be simply estimated with an empirical distribution (plug-in), we focus on the error induced by approximating $P_{Y, A | X}$ using an estimate 
    $\what{P}_{Y, A | X}$.
    \begin{align*}
        P_X[\what{\mu}(X)] - P_X [\mu(X)]
        & =     P_X [\what{P}[Y \mid A = 1, X]]
        - P_X [P[Y \mid A = 1, X]] \\
        & = \psi(P_X \times \what{P}_{Y, A | X}) - 
    \psi(P_X \times P_{Y, A | X}).
    \end{align*}
    \looseness=-1 For the statistical functional $Q \mapsto \psi(Q)$, let $\varphi$ be its canonical gradient (a.k.a. efficient influence function) 
    with respect to $Q_{Y, A|X}$; 
    without loss of generality, we require 
    $Q[\varphi(Z; Q) | X] =0$ for all $Q$. 
    Then, a (distributional) first-order Taylor expansion gives
    \begin{align}
        \psi(P_X \times \what{P}_{Y, A | X}) 
        -  \psi(P_X \times P_{Y, A | X})
        & = \iint \varphi(Z; \what{P}) d (\what{P}_{Y, A|X} - P_{Y, A|X}) dP_X + R_2(\what{P}, P) \nonumber \\
        & = - \iint \varphi(Z; \what{P}) dP_{Y, A|X} dP_X + R_2(\what{P}, P) \nonumber \\
        & = - P [ \varphi(Z; \what{P}) ] 
        + R_2(\what{P}, P)
        \label{eqn:taylor}
    \end{align}
    where $R_2$ is a second-order remainder term.
    Conclude that the first-order error term of the plug-in approach $P_X[\what{\mu}(X)]$ is given by $-P [ \varphi(Z; \what{P}) ]$.
\fi

A common approach in semiparametric statistics to better estimate $\psi(P)$ is to explicitly correct for this first-order error term. 
As we discuss later in \cref{sec:theory_text}, the first-order correction leads to asymptotic optimality properties like semiparametric efficiency---providing the shortest possible confidence interval---and double robustness---achieving estimator consistency even if only one of $\what\mu(X), \what\pi(X)$ is consistent. 

We let $\varphi_X$ denote a \emph{projected} version of the canonical gradient, 
where we keep only the component that remains when the distribution of $X$ is fixed. Practically speaking, the methods below that empirically correct for this projected version produce the same guarantees as empirically correcting for the full canonical gradient \cite{VanDerLaanRo11,van2006targeted}.\ifdefined\msom\else\ifdefined\pnas\yf{footnote with broken reference}\footnote{We also discuss this briefly in \cref{sec:projected_if}.}. 
\else\fi\fi
\paragraph{First-Order Correction for the ATE} For the ATE~\eqref{eqn:ate}, it is well-known that 
\begin{equation}
\label{eqn:ate-pathwise}
\varphi_X(Z; P)
= \frac{A}{P[A =1 | X]} (Y - P[Y | A = 1, X])
= \frac{A}{\pi(X)} (Y - \mu(X))
\end{equation}
which we do not derive in this work. %
For an accessible primer on such derivations, see~\citep{Kennedy22}. 

\paragraph{One-Step Estimation (Augmented Inverse Propensity Weighting, AIPW)~\cite{Bang2005DoublyRE}} 
One of the most prominent first-order correction approaches modifies the plug-in estimator by moving the first-order error term to the left-hand side in the Taylor expansion~\eqref{eqn:taylor}. The resulting estimator achieves second-order error rates:
\vspace{-0.3em}
\begin{align*}
    P\left[ \what{\mu}(X)
    + \varphi_X(Z; \what{P})\right]
    - P[\mu(X)] = R_2(\what{P}, P).
\end{align*}
Using the empirical distribution with samples $Z_1,\ldots,Z_N$ to approximate $P$ in 
the one-step debiased estimator, we arrive at the augmented inverse propensity weighted estimator
\begin{align}
    \what\psi^{\textrm{\rm AIPW}} 
     & := \frac{1}{n} \sum_{i=1}^n 
    \left(\what{\mu}(X_i)
    + \varphi_X(Z_i; \what{P}) \right) %
     = \frac{1}{n} \sum_{i=1}^n 
    \what{\mu}(X_i)
    + \frac{1}{n} \sum_{i=1}^n \frac{A_i}{\what{\pi}(X_i)} 
    (Y_i - \what{\mu}(X_i)) .
        \label{eqn:debiased}
\end{align}

\paragraph{Targeting (Targeted Maximum Likelihood Estimation)~\cite{VanDerLaanRo11,van2006targeted}} 
Targeting takes an alternative and more general approach to first-order correction. It commits to the use of the plug-in estimator, and constructs a tailored adjustment (``fluctuation'') to the existing 
nuisance parameter estimate %
to set the 
first-order error to zero, where the magnitude of the adjustment is solved by maximizing likelihood (hence the name). The fluctuation takes place in the 
direction of a task-specific random variable (``clever covariate''),  which takes 
different forms depending on the estimand. For illustrative purposes, in this discussion, we use focus on a formulation of targeting for unbounded outcomes under the squared loss function. %
In our experiments, we also include a commonly used variant of targeting that increases the stability of the estimator by assuming bounded outcomes, and modeling the outcomes with a logistic link.
\footnote{The literature on targeting is large, with variants including regularization techniques and adaptive versions \citep{VanDerLaanRo11,van2011cross,van2006targeted,van2024adaptive}. 
We choose to focus on the formulation in Equation~\eqref{eqn:targeting} for its simplicity. See %
the discussion on TMLE for bounded outcomes in \cref{sec:ks}, and \cref{sec:tmle_comparison} for additional discussion.
}
We defer a deeper discussion of TMLE to \Cref{sec:tmle_comparison}.

We begin by using standard notation for TMLE, which has notation for both $A=0$ and $A=1$. We apply it to our setting with $Y(0):=0$ (so that $\what\mu(X)\defeq \what\mu(1,X)$, as in \Cref{sec:ate}) and consider general outcomes $Y$
(including unbounded continuous $Y$). The form of the canonical gradient~\eqref{eqn:ate-pathwise} motivates the use of $H(A,X):=\frac{A}{ \what{\pi}(X)}$ as the clever covariate: %
targeting uses an adjusted nuisance estimate
$\what{\mu}(A,X)+\epsilon^\star H(A,X)$ %
in place of $\what{\mu}(A,X)$
in the plug-in for $\psi(P)=P[\mu(1,X)]$, 
where $\epsilon^\star$ is chosen to solve the targeted maximum likelihood problem
\vspace{-0.3em}
\begin{equation}
    \label{eqn:targeting}
    \epsilon^\star:=\argmin_{\epsilon \in \R} 
    ~~\frac{1}{n}\sum_{i=1}^n  A_i \left( 
    Y_i - \what{\mu}(A_i,X_i) - \epsilon \frac{A_i}{\what\pi(X_i)}\right)^2.
\end{equation}
From the KKT conditions, %
the solution
removes the 
finite-sample estimate of the first-order error term~\eqref{eqn:taylor} of the plug-in estimator using $\what{\mu}(A,X) + \epsilon^\star \frac{A}{ \what{\pi}(X)}$. Solving for $\epsilon\opt$, we obtain
\vspace{-0.3em}
\begin{equation}
    \label{eqn:tmle-eps}
    \epsilon^\star = 
    \left(
    \frac{1}{n}\sum_{i=1}^n \frac{A_i}{\what{\pi}(X_i)^2}
    \right)^{-1}
    \frac{1}{n} \sum_{i=1}^n 
    \frac{A_i}{\what{\pi}(X_i)} (Y_i - \what{\mu}(X_i)).
\end{equation}
Thus, we arrive at an explicit formula for the targeted maximum likelihood estimator
\vspace{-0.3em}
\begin{align}
    \what{\psi}^{\rm TMLE}
    & \defeq \frac{1}{n} \sum_{i=1}^n \Big(\what\mu(1,X_i)+\epsilon^\star H(1,X_i)\Big) %
    =
    \frac{1}{n} \sum_{i=1}^n \left(\what{\mu}(X_i)
    + \epsilon^\star \frac{1}{\what{\pi}(X_i)}\right) \nonumber \\
    & = \frac{1}{n} \sum_{i=1}^n \what{\mu}(X_i)
    + \frac{\sum_{i=1}^n \frac{1}{\what{\pi}(X_i)}}{\sum_{i=1}^n \frac{A_i}{\what{\pi}(X_i)^2}}
    \cdot \frac{1}{n} \sum_{i=1}^n 
    \frac{A_i}{\what{\pi}(X_i)} (Y_i - \what{\mu}(X_i)).
        \label{eqn:tmle}   
\end{align}



















\paragraph{Targeted Regularization~\cite{chernozhukov2022riesznet,shi2019adapting}}
Instead of modifying the outcome model in the direction of the clever covariate in a post-processing step~\eqref{eqn:targeting},
we can instead regularize the usual training objective for the outcome model using a similar adjustment term throughout training, an approach referred to as \emph{targeted regularization}, %
where it is demonstrated on e.g. neural networks learned via stochastic gradient-based optimization. %
\ifdefined\msom\else\footnote{An additional difference between TMLE and targeted regularization is the data splits involved. In TMLE, the targeting objective~\eqref{eqn:targeting} for calculating $\epsilon^*$ uses the same data that is used for plug-in estimation (the first line in \eqref{eqn:tmle}),while in targeted regularization, the targeting objective is applied to the dataset used for training $\what\mu$, which has not explicitly described above, as so far, $\what\mu, \what\pi$ are taken as given. These data splits are often different; see \cref{sec:data_splitting} for more discussion on data splits.}\fi %
